\section{Código abierto, direcciones futuras y preguntas y respuestas de la clase}

Al cierre de la clase, NLLB no se presenta como un punto final: se proponen tres direcciones futuras y, en la sesión de Q\&A, se añaden precisiones sobre la participación de la comunidad, el volumen del trabajo con datos, Common Crawl, zero-shot, speech translation y los límites de los foundation models. Este contenido oral explica cómo el sistema especializado NLLB de 2022 se conecta con los modelos multimodales de propósito general que llegaron después.

\subsection{Explicit multilinguality, support y facilidad de uso}

La slide de direcciones futuras destaca explicit multilinguality, support for everyone, y que el modelo sea más fácil de usar y de entrenar. El primer punto exige que el modelo conozca de forma explícita el idioma, la escritura y la dirección de traducción; el segundo incluye a los idiomas de la cola entre los objetivos; el tercero reconoce que una cadena de procesamiento de datos que solo pueden operar las grandes instituciones todavía no basta para constituir una infraestructura de acceso universal.

\lecturefigure{slide-39-future-directions.jpg}{Direcciones futuras de NLLB: multilingüismo explícito, cobertura más amplia y menores barreras de uso}{Video de la clase de Stanford 00:41:15}

\begin{knowledgebox}{Lectura de la figura: las direcciones futuras corresponden a tres niveles, representación, cobertura e ingeniería}
Explicit multilinguality actúa sobre los language tags, el routing y la evaluation; support for everyone actúa sobre los idiomas de la cola, las variantes dialectales y las necesidades de las comunidades; ease of use/training actúa sobre las herramientas abiertas, la barrera de cómputo y la adaptación local. Si falta cualquiera de los tres, los resultados de investigación difícilmente se convierten en una capacidad pública utilizable de forma sostenida.
\end{knowledgebox}

\teachervoice{Fan señala en particular que algunos idiomas son sobre todo orales y no escritos. Un NLLB de texto no puede cubrir esos idiomas a partir de nada; en el futuro habrá que conectar speech, transcription y text translation, en lugar de tratar la «ausencia de datos web» como si el idioma no existiera.}

\subsection{El trabajo con datos y el trabajo con modelos pesan casi lo mismo}

En la sesión de Q\&A alguien preguntó cómo se repartían los recursos del proyecto. Fan respondió con una estimación aproximada: cerca de la mitad fue trabajo de data o data-adjacent, y la otra mitad de modeling, aunque la frontera es muy difusa. LID, el encoder, el filtering y la evaluation pertenecen a la vez a los datos y al modelo; lo importante de esta respuesta no es una proporción exacta de gestión de proyecto, sino refutar la idea de que «el trabajo principal es entrenar un modelo grande».

\begin{importantbox}{Niveles de los activos de ingeniería de NLLB}
\begin{tabularx}{\textwidth}{@{}L{0.2\textwidth}L{0.29\textwidth}X@{}}
\toprule
Nivel & Activos públicos & Pregunta que permiten reproducir \\
\midrule
Need / standards & Conclusiones de las entrevistas y proceso de normalización lingüística & Por qué se definieron así los objetivos y las etiquetas \\
Evaluation & FLORES-200, human protocol, Toxicity-200 & Si el modelo mejora y dónde no es seguro \\
Data tools & LID, LASER3, Stopes, scripts de reconstrucción & Cómo se generan y filtran los datos de entrenamiento \\
Models & Checkpoints y código de NLLB & Comportamiento final en inferencia y adaptación \\
\bottomrule
\end{tabularx}
Una apertura completa significa cubrir, en la medida de lo posible, toda la cadena que va del problema al modelo, y no solo publicar el checkpoint.
\end{importantbox}

\teachervoice{Al decir «50-50», el docente añadió de inmediato que data y modeling difícilmente se separan con rigor. Estas notas lo toman como una intuición de orden de magnitud: en los grandes proyectos multilingües, una parte considerable de la innovación principal ocurre en los datos, la evaluación y la cadena de herramientas.}

\subsection{Que Common Crawl sea abierto no significa que represente por igual}

En clase se explicó que Common Crawl es un conjunto de instantáneas web abiertas y descargables, y se usó una formulación prudente sobre su frecuencia de actualización. El acceso abierto resuelve parte del problema de los permisos para obtener datos, pero no resuelve la proporción de páginas web por idioma, la duplicación, las plantillas, el contenido basura, el sesgo geográfico ni la deriva temporal.

\begin{warningbox}{Web-scale y language coverage son dos variables distintas}
Un corpus puede contener miles de millones de páginas web y, al mismo tiempo, casi no tener texto útil en ciertos idiomas. Al evaluar un mining pipeline conviene reportar, por idioma, las raw pages, la LID pass rate, los tokens tras la deduplicación, el bitext candidato, la precision humana y la distribución por dominio, en lugar de informar solo el tamaño total del crawl.
\end{warningbox}

\subsection{Los límites de zero-shot: una dirección no vista no es un idioma no visto}

NLLB no extrajo datos de entrenamiento para cada dirección $A\rightarrow B$; la explosión combinatoria lo haría costoso y no necesariamente acorde con las necesidades reales. Una dirección puede ser zero-shot, pero si el modelo ya vio los idiomas $A$ y $B$ en otras direcciones, todavía puede transferir conocimiento mediante las representaciones compartidas del encoder/decoder.

\begin{importantbox}{Hay que distinguir dos tipos de «no visto»}
\begin{itemize}
  \item \textbf{Unseen direction}: el entrenamiento no incluyó directamente $A\rightarrow B$, pero sí tareas como $A\leftrightarrow E$ y $B\leftrightarrow E$.
  \item \textbf{Unseen language}: el texto, los tokens o la supervisión del idioma $A$ casi nunca aparecieron.
\end{itemize}
En clase, zero-shot se refiere sobre todo al primer caso. Presentarlo como capacidad de traducir un idioma completamente nuevo exagera el alcance de la generalización.
\end{importantbox}

\teachervoice{Fan explica que algunos zero-shot pairs funcionan muy bien porque el modelo ya vio el idioma de entrada y el de salida, aunque no esa combinación dirigida; en conjunto, estas direcciones siguen siendo más débiles, y no se trata de traducción sin ejemplos previos de un unseen language por completo.}

\subsection{Del NLLB de texto a SeamlessM4T}

Ante la pregunta sobre voz, Fan aclaró que NLLB en sí es principalmente traducción de texto; fue el trabajo posterior, SeamlessM4T, el que integró speech, transcription y text translation. Esta respuesta separa con claridad los límites del proyecto: los supuestos sobre datos del pipeline de texto no pueden extrapolarse directamente a idiomas que dependen sobre todo de la lengua oral.

\begin{knowledgebox}{Cadena de evidencia de la extensión de texto a voz}
Speech translation requiere audio, transcripciones, cobertura de hablantes y acentos, segmentación del habla y representaciones acústicas; el NLLB de texto trabaja sobre todo con pares de oraciones escritas. SeamlessM4T conecta ambos mediante representaciones multimodales compartidas, pero sigue necesitando nuevas evaluaciones de voz y nuevos análisis de seguridad; no basta con reutilizar las puntuaciones de texto de FLORES.
\end{knowledgebox}

\subsection{Una nueva lectura en la era de los foundation models}

Al final de la clase se discutió que, si el proyecto se reiniciara en una era de modelos base más potentes, podría usarse un foundation model grande con supervised fine-tuning, lo que reduciría la necesidad de diseñar sistemas especializados para cada tarea. Fan recordó que en los inicios del NLP la investigación en translation, summarization y QA estaba separada, mientras que ahora cada vez más tareas se apoyan en una base compartida.

\teachervoice{El juicio del docente es una tendencia, no un teorema: cuanto más potente sea el modelo base, menos fine-tuning de dominio podría requerirse, pero sigue siendo una condición previa que los idiomas de bajos recursos estén bien representados en el preentrenamiento. Una base de propósito general no corrige por sí sola el desequilibrio de idiomas en el corpus de entrenamiento.}

\begin{warningbox}{No hay que borrar la contribución de NLLB con el relato posterior de los LLM}
Un foundation model más potente sigue necesitando datos multilingües, benchmarks confiables, language identification, evaluación humana y safety probes. El valor de largo plazo de NLLB está justamente en haber publicado de forma conjunta esa infraestructura y los modelos, no solo en haber propuesto una arquitectura de 2022.
\end{warningbox}

\subsection{Resumen de la sección}

Los sistemas multilingües futuros deben abarcar una representación explícita del idioma, las comunidades de la cola y barreras de ingeniería más bajas. La sesión de Q\&A añade precisiones: los datos y el modelo son igual de decisivos, la web abierta no está equilibrada, zero-shot se refiere sobre todo a direcciones no vistas, speech requiere evidencia de una nueva modalidad y un foundation model tampoco sustituye la infraestructura multilingüe.
